\documentclass{ximera}
\title{Examples}

\begin{document}

    \begin{abstract}   \end{abstract}
    \maketitle

\begin{example}
$\mathbb{R}^n$ is a smooth manifold in a canonical way. %That is, there is a unique smooth structure such that for each open set $U \subset \mathbb{R}^n$, the inclusion $U \hookrightarrow \mathbb{R}^n$ is a chart.  
\end{example}

\begin{proof}
    $\mathbb{R}^n$ is a metric space,  hence Hausdorff.   The set of balls with rational radius and center with rational coordinates forms a countable sub-basis of the topology of $\mathbb{R}^n$, so it is second countable. 
    
    The identity map $\text{Id}_{\mathbb{R}^n} : \mathbb{R}^n \to \mathbb{R}^n$ is a chart. Since the domain is everything, and any chart is compatible with itself, the set containing this chart is a smooth atlas. More specifically, $\mathcal{A} : = \{ (\mathbb{R}^n, \text{Id}_{\mathbb{R}^n} ) \}$ is a smooth atlas of the topological manifold $\mathbb{R}^n$. 
    
    Therefore, there is a unique smooth structure containing this chart. Note that for each open set $U \subset \mathbb{R}^n$, the inclusion $U \hookrightarrow \mathbb{R}^n$ belongs to this smooth structure.
\end{proof}



\begin{example}[Open sets]
Let $M$ be an $n$-dimensional smooth manifold and $U \subset M$ an open set. Then $U$ is an $n$-dimensional smooth manifold. 
\end{example}

\begin{proof}[Proof sketch]
    Let $\mathcal{S} = \{ (U_i, \varphi _i ) \} _{i \in I }$ be the smooth structure of $M$. Then $\mathcal{S}_U : = \{ ( U \cap U_i , \varphi_i \vert _{U \cap U_i}  ) \} _{i \in I }$ is a smooth atlas on $U$ (it is actually a smooth structure). 
\end{proof}


\begin{example}[Locally a graph] Let $M\subset \mathbb{R}^{n+m}$ be a set that is locally the graph of a smooth function. That is, for each $p \in M$, there is a linear isomorphism $L : \mathbb{R}^{n+m} \to \mathbb{R}^{n+m}$, open sets $U \subset \mathbb{R}^n $, $V \subset \mathbb{R}^m$, and a smooth function $f : U \to \mathbb{R}^n$ such that $L(p) \in U\times V$ and
\[      L (M) \cap ( U \times V  ) = \{ (x,f(x)) \in \mathbb{R}^n \times \mathbb{R}^m  \, \vert \,  x \in U \}  . \]
Then $M$ is a smooth manifold. 
\end{example}



\begin{proof}
    Let $L : \mathbb{R}$, $U$, $V$, and $f$ as above. Define
    \[ \varphi  : M  \cap L^{-1} (U \times V) \to U \]
    as 
    \[  \varphi : = \pi \circ L             \]
    where $\pi  : \mathbb{R}^n \times \mathbb{R}^m \to \mathbb{R}^n$ denotes the projection onto the first $n$ coordinates.  Notice that  
    \begin{align*}
        \varphi (L^{-1} (x, f(x))) & = \pi \circ L \circ L^{-1} (x, f(x)) \\
        & = \pi (x, f(x)) \\
        & = x .
    \end{align*}
     Hence the inverse $\varphi^{-1} : U \to M \cap L^{-1} (U \times V)$ is the continuous function  
    \[ \varphi  ^{-1} (x) =  L ^{-1} (x, f(x)) ,  \]
    showing that $\varphi$ is a chart. We claim that the set of all charts obtained this way form an atlas. For that purpose, take a linear isomorphism $ \hat{L} : \mathbb{R}^{n+m} \to \mathbb{R}^{n+m}$, open sets $\hat{U} \subset \mathbb{R}^n $, $\hat{V} \subset \mathbb{R}^m$, and a smooth function $\hat{f} : \hat{U} \to \mathbb{R}^n$ such that 
    \[      \hat{L} (M) \cap ( \hat{U} \times \hat{ V}  ) = \{ (x,\hat{f}(x)) \in \mathbb{R}^n \times \mathbb{R}^m  \, \vert \,  x \in \hat{U} \}  . \]
    Construct with them the chart 
    \[             \psi : M \times \hat{L}^{-1} (\hat{U}  \times \hat{V} ) \to \hat{U}                            \]
    given by 
    \[     \psi : = \pi \circ \hat{L} .     \]
    The change of coordinates is then given by
    \[         \psi \circ \varphi ^{-1} (x) = ( \pi \circ \hat{L} ) \, ( L ^{-1} (  x, f( x )  ) )  ,                                     \]
    which is the composition of the smooth functions
    \[       x \mapsto ( x, f(x) ) \mapsto L ^{-1} (x, f(x)) \mapsto \hat{L} ( L ^{-1} (x,f(x)) ) \mapsto \pi (  \hat{L}  ( L^{-1}  (x, f(x))  )  ).                            \]
    Therefore the charts $\varphi$ and $\psi $ are compatible, proving our claim. 
    \end{proof}


    \begin{example}[Level sets]
        Let $\Omega \subset \mathbb{R}^{n+m} $ be an open set and $F \in C^{\infty} ( \Omega ; \mathbb{R}^m  )$. If for all $p \in F^{-1} (0)$, the differential $d_pF : \mathbb{R}^{n+m} \to \mathbb{R}^m$ is surjective, then the level set $F^{-1}(0) \subset \mathbb{R}^{n+m}$ is a smooth manifold.  
    \end{example}


    \begin{proof}
    By the Implicit Function Theorem, this is covered by the  above example
    \end{proof}



    \begin{example}[Products]
        Let $M$ and $N$ be smooth manifolds of dimension $m$ and $n$, respectively. Then $M \times N$ is a smooth manifold of dimension $m + n$ in a canonical way.
    \end{example}

    \begin{proof}
        For each chart $(U, \varphi)$ of $M$ and each chart $(V, \psi) $ of $N$, we consider the map
        \[    \varphi \otimes \psi : U \times V \to \varphi (U) \times \psi (V)  \subset \mathbb{R}^{n+m}       \]
        given by 
        \[  (\varphi \otimes \psi ) (x,y) : = (\varphi (x), \psi (y)). \]
        Since products of homeomorphisms are homeomorphisms, $\varphi \otimes \psi$ is a chart.  
        
        
        Now consider  $(U', \varphi ' )$ another chart of $M$ and $(V' , \psi ' ) $ another chart of $N$. With those, we build the chart
        \[      \varphi ' \otimes \psi ' : U ' \times V' \to \varphi ' (U ' ) \times \psi ' (V ' )   .        \] 
        Notice that 
        \[    (U \times V) \cap (U ' \times V' ) = (U \cap U') \times (V \times V')   .       \]
        Then the change of coordinates between $\varphi \otimes \psi $ and $\varphi ' \otimes \psi ' $ is the function 
        \[       (\varphi ' \otimes \psi ') \circ (\varphi \otimes  \psi ) ^{-1} : ( \varphi  \otimes \psi ) ( ( U \cap U ' ) \times (V \cap V') )  \to ( \varphi ' \otimes \psi ' ) ( ( U \cap U ' ) \times (V \cap V') )                    \]
        given by 
        \begin{align*}
             (\varphi ' \otimes \psi '  ) \circ (\varphi \otimes \psi ) ^{-1}  (x,y)  & = (\varphi ' \otimes \psi '  ) (\varphi ^{-1} (x) , \psi ^{-1}   y)  \\
             & =  ( \varphi ' \circ \varphi ^{-1} (x) , \psi ' \circ \psi ^{-1} (y) ) .
        \end{align*}
        Since the maps $\varphi ' \circ \varphi ^{-1}$ and $\psi ' \circ \psi ^{-1}$ are smooth, the change of coordinates is smooth.
    \end{proof}



\begin{example}
    Let $M_n (\mathbb{R})$ denote the space of $n\times n$ real matrices and identify it with $\mathbb{R}^{n^2}$. Given $A \in M_n(\mathbb{R})$, we denote by $A^T $ its transpose. Then 
    \[      O  (n ) : = \{ A \in M_n (\mathbb{R}) \, \vert \, A A^{T}  = 1 \}         \]
    is a smooth manifold of dimension $n(n-1)/2$. 
\end{example}


\begin{proof}[Proof]
        We identify $M_n (\mathbb{R}) = \mathbb{R}^{n^2}$ with $(\mathbb{R}^n)^n$, the set of $n$-tuples of elements of $\mathbb{R}^n$. Under this identification, $O (n)$ consists of the set of orthonormal bases.  Define 
        \[    F :    (\mathbb{R}^n)^n \to \mathbb{R}^{n(n+1)/2}      \]
        by
        \[      F _{ij}(  \Vec{v}_1, \ldots,  \Vec{v}_n ) = \Vec{v}_i \cdot \Vec{v}_j - \delta _{ij}                 \]
        for $1\leq i \leq j\leq n$, where $\delta_{ij}$ denotes the Kronecker delta. Then
        \[        F^{-1} (0) = O (n) .                       \]
        Fix $  p : =  (\Vec{v}_1, \ldots , \Vec{v}_n) \in O (n) ,              $        and $1 \leq i_0 \leq j_0 \leq n$. Consider the curve 
        \[    \gamma ( t )  : = (  \Vec{v}_1 ,  \ldots , \Vec{v}_{j_0} + t \Vec{v}_{i_0} , \ldots , \Vec{v}_n )        .            \]
        For $1 \leq i \leq j \leq n$, we have
        \begin{align*}
           F _{ij } (\gamma (t))  & =  ( \Vec{v}_i + t \,\delta_{ij_0}  \Vec{v}_{i_0} ) \cdot ( \Vec{v}_j + t \, \delta_{jj_0} \Vec{v}_{i_0} ) \\                   & = \Vec{v}_i \cdot \Vec{v}_j + t ( \delta_{ij_0} \Vec{v}_{i_0} \cdot \Vec{v}_j + t \delta_{jj_0} \Vec{v}_i \cdot \Vec{v}_{i_0} ) + t^2 \delta_{ij_0} \delta_{jj_0} \\
           & = \delta_{ij} + t (\delta_{ij_0}\delta_{i_0j} + \delta_{jj_0}\delta_{ii_0}) + t^2 \delta_{ij_0} \delta_{jj_0}. 
        \end{align*}
        Since $i \leq j$ and $i_0 \leq j_0$, we have
        \[      \delta_{ij_0}\delta_{i_0j} + \delta_{jj_0}\delta_{ii_0} = \begin{cases}
            2 & \text{ if } i=i_0 =  j= j_0 \\
            1 & \text{ if } i=i_0 < j= j_0 \\
            0 & \text{ otherwise}
        \end{cases}        .            \]
        Therefore, 
        \[  (F \circ \gamma )'_{ij} (0) = \begin{cases}
            2 & \text{ if } i=i_0 =  j= j_0 \\
            1 & \text{ if } i=i_0 < j= j_0 \\
            0 & \text{ otherwise}
        \end{cases}              .                    \]
        Since $\gamma (0) = p$, by the chain rule we have 
        \[  d_p F (\gamma' (0)) _{ij} = (F \circ \gamma ) '_{ij} (0)  =   \begin{cases}
            2 & \text{ if } i=i_0 =  j= j_0 \\
            1 & \text{ if } i=i_0 < j= j_0 \\
            0 & \text{ otherwise}
        \end{cases}   .                           \]
        This shows that the vector in $\mathbb{R}^{n(n+1)/2}$ that is $1$ in the $(i_0, j_0)$ coordinate and $0$ in all other coordinates is in the image of $d_pF$. Therefore $d_pF $ is surjective. Since $p$ was arbitrary, $O (n) $ is a smooth manifold of dimension $n(n-1)/2$.
        
        \end{proof}







\end{document}